\documentclass[11pt]{article}

\usepackage[margin=1in]{geometry}
\usepackage{amsmath,amssymb,amsthm,mathtools,booktabs,tabularx}
\usepackage{microtype,xurl}
\usepackage{xcolor}
\usepackage[colorlinks=true,linkcolor=blue!55!black,
            citecolor=blue!55!black,urlcolor=blue!55!black]{hyperref}

\newtheorem{theorem}{Theorem}[section]
\newtheorem{proposition}[theorem]{Proposition}
\newtheorem{corollary}[theorem]{Corollary}
\theoremstyle{definition}
\newtheorem{definition}[theorem]{Definition}
\theoremstyle{remark}
\newtheorem{remark}[theorem]{Remark}

\newcommand{\R}{\mathbb R}
\newcommand{\C}{\mathbb C}
\newcommand{\dd}{\mathrm d}
\newcommand{\Lie}{\mathcal L}
\newcommand{\Ric}{\operatorname{Ric}}
\newcommand{\Diff}{\operatorname{Diff}}
\newcommand{\Weyl}{\operatorname{Weyl}}
\newcommand{\cert}[1]{\nolinkurl{#1}}
\newcolumntype{Y}{>{\raggedright\arraybackslash}X}

\title{Fourth-Order Radiation of Pure-Weyl Black Holes\\
\large Exact Branch Splitting, Horizon Flux, and a Causal Truncation No-Go}
\author{GPT-5.6.sol\thanks{OpenAI model and principal programme author.}
\and Claude Fable 5\thanks{Anthropic model and computational coauthor. The
project was commissioned and coordinated by Asger Alstrup Palm
(\texttt{asger@area9.dk}), who initiated the questions, exercised editorial
control, and serves as corresponding human contact, but claims no technical
authorship.}}
\date{Working draft, 19 July 2026}

\begin{document}
\maketitle

\begin{abstract}
We test whether the Einstein radiative branch of four-dimensional pure Weyl
gravity can be selected by ordinary causal boundary conditions around a black
hole.  We first classify an exact static spherical Bach-flat family within a
declared Laurent ansatz, reproduce the Mannheim--Kazanas solution, and derive
a residual-basic horizon generator.  On the static parameter slice its
Lee--Wald Hamiltonian, Wald entropy, and temperature satisfy an exact first
law at every simple horizon.  Arbitrary time-dependent spherical Weyl and
diffeomorphism perturbations are null charge directions at linear order.

We then study Schwarzschild perturbations.  In the axial \(\ell=2\) sector
the fourth-order Bach equation factors exactly through the linearized Ricci
tensor:
\[
 \delta B_{ab}=\frac12\Box(\delta R_{ab})
                +C_{acbd}(\delta R)^{cd}.
\]
The Einstein/Regge--Wheeler branch is \(\delta R_{ab}=0\); the extra branch
is a second-order Lichnerowicz-type wave for
\(\psi_{ab}=\delta R_{ab}\).  The extra branch has a two-parameter analytic
ingoing family at the future horizon, is luminal and bounded at infinity for
real frequency, and has the same leading asymptotic characteristics as the
Einstein branch.  The action-derived Lee--Wald current makes the Einstein
self-block null for conjugate waves, while independently verified fixtures
give nonzero Einstein--extra and extra--extra horizon flux.

It follows that horizon regularity plus local causal decay cannot truncate
the axial fourth-order exterior to its Einstein branch.  Such a truncation
would require an additional projection on scattering data rather than a
local initial-boundary condition.  In the polar \(\ell=2\) sector we prove
the corresponding exact trace-coupled branch split; its horizon reach, flux,
and causal disposition remain the next theorem gate.  We make no claim about
complex frequencies, stability, quasinormal spectra, nonlinear evolution,
Hawking radiation, or quantum unitarity.
\end{abstract}

\section*{AI authorship and computational accountability}
The named AI systems developed the programme, derivations, symbolic
implementations, independent verification architecture, claim audits, and
this manuscript.  Asger Alstrup Palm commissioned and coordinated the work
and is the corresponding human contact.  Every certified statement below is
bound to a machine-readable exact certificate and a separately implemented
verifier.  The cross-flux result uses controlled high-precision series
evaluation with exact recurrence data and independent frequencies; all local
tensor identities and polynomial classifications are exact.  The manuscript
remains conditional on final human review.

\tableofcontents

\section{The physical question}

Fourth-order gravitational equations contain more local solutions than the
Einstein equations.  A common proposal is to retain the Einstein branch by a
boundary condition.  Around a Lorentzian black hole, that proposal must pass
a causal test:

\begin{quote}
Can regular initial data and local conditions at the future horizon and
outer boundary remove the extra branch while retaining ordinary Einstein
gravitational waves?
\end{quote}

The answer in the certified Schwarzschild axial \(\ell=2\) laboratory is
no.  The extra branch is regular at the future horizon, propagates on the
same real characteristics as the Einstein branch, has no growing real-
frequency asymptotic solution, and carries nonzero action-derived horizon
flux.  Ordinary local regularity and decay therefore admit it.

This conclusion is narrower than a statement about the full theory.  It is
an axial, linear, real-frequency causal-disposition theorem.  The polar
branch has been split but not yet propagated through the same chain, and no
complex-frequency stability analysis has been performed.

\subsection{Claims and dependency tags}

The static family and local operator identities carry
\textsc{local-algebraic}.  Harmonic branch and current statements also carry
\textsc{reduced-mode}.  The final axial disposition uses the declared
horizon and asymptotic domains and is a scoped \textsc{lorentzian-causal}
no-go.  It is not a complete exterior Green-hyperbolicity theorem.

\section{Pure Weyl gravity and its covariant current}

We use signature \((-+++)\) and action
\begin{equation}
 S_W[g]=\alpha\int_M\sqrt{-g}\,C_{abcd}C^{abcd}\,\dd^4x.
\end{equation}
The field equation is the Bach equation \(B_{ab}=0\), with convention
\begin{equation}
 B_{ab}=\nabla^c\nabla^dC_{acbd}
       +\frac12R^{cd}C_{acbd}.
\end{equation}
The covariant presymplectic current is obtained from the literal variation
of the curvature-momentum potential of this action.  Its normalization is
fixed by the off-shell Green identity
\begin{equation}\label{eq:green-identity}
 \partial_tF^t(h_A,h_B)+\partial_rF^r(h_A,h_B)
 =4\alpha\int_{S^2}\sqrt{q}\,
 \bigl(h_B^{ab}\delta B_{ab}[h_A]
      -h_A^{ab}\delta B_{ab}[h_B]\bigr).
\end{equation}
Thus the current is conserved for pairs of linearized Bach solutions and is
not a freely rescaled master-equation Wronskian.

\section{Static spherical Bach-flat family}

In Schwarzschild gauge,
\begin{equation}
 \dd s^2=-B(r)\dd t^2+B(r)^{-1}\dd r^2+r^2\dd\Omega_2^2,
\end{equation}
the declared Laurent ansatz is
\begin{equation}
 B(r)=w-\frac{u}{r}+\gamma r-kr^2+\frac{c_2}{r^2}+c_3r^3.
\end{equation}

\begin{theorem}[Complete Bach-flat locus in the Laurent ansatz]
The metric is Bach-flat if and only if
\begin{equation}
 c_2=c_3=0,
 \qquad
 w^2+3u\gamma=1.
\end{equation}
The Mannheim--Kazanas parametrization is recovered by
\[
 w=1-3\beta\gamma,qquad
 u=\beta(2-3\beta\gamma).
\]
Within this family the metric is Einstein if and only if \(\gamma=0\).
\end{theorem}

Completeness here refers only to the Laurent ansatz.  The certificate does
not claim a new proof of the global static spherical classification.

The form-preserving local \(\Diff\times\mathrm{Weyl}\) transformations act
as dilation and translation on the cubic
\[
 Q(x)=-ux^3+wx^2+\gamma x-k,qquad x=r^{-1}.
\]
Their orbit rank is two.  A single continuous invariant may be taken as
\[
 J=u^2\operatorname{disc}(Q),
\]
while \(\operatorname{disc}(Q)=0\) is the degenerate-horizon locus.
The conformal translation may fail globally where its Weyl factor vanishes;
it is used only as a local orbit statement.

An exact non-Einstein three-horizon fixture is
\[
 (\beta,\gamma,k)=\left(\frac32,\frac{12}{19},\frac1{19}\right),
 \qquad
 rB(r)=-\frac1{19}(r-1)(r-3)(r-8).
\]
All curvature invariants are finite at the simple horizons.  In ingoing
Eddington--Finkelstein coordinates,
\[
 \dd s^2=-B(r)\dd v^2+2\dd v\dd r+r^2\dd\Omega_2^2,
\]
the metric and Bach equation are regular there.

\section{Residual-basic energy and the static first law}

For the chart generator \(\partial_t\), the bare static Lee--Wald one-form
is not integrable.  The obstruction disappears when the generator is
normalized consistently with the residual quotient.

\begin{theorem}[Normalized static generator]
Let
\begin{equation}
 \chi=u\,\partial_t,
 \qquad u=\beta(2-3\beta\gamma).
\end{equation}
Up to a function of the single residual invariant \(J\), this normalization
is forced by residual basicness and dilation weight.  Its Hamiltonian is
\begin{equation}
 H=-16\pi\alpha\beta^2D_2,
 \qquad
 D_2=9\beta^2\gamma^2k-\beta\gamma^3
     -12\beta\gamma k+\gamma^2+4k.
\end{equation}
At every simple horizon \(r_h\), the Wald entropy
\begin{equation}
 S_h=64\pi^2\alpha\beta
 \frac{2-3\beta\gamma+\gamma r_h}{r_h}
\end{equation}
and normalized temperature
\begin{equation}
 T_h=\frac{uB'(r_h)}{4\pi}
\end{equation}
satisfy
\begin{equation}
 \dd H=T_h\,\dd S_h
\end{equation}
identically modulo \(B(r_h)=0\).
\end{theorem}

This is a theorem on the static parameter slice, not yet a general
physical-process first law.  Its linear dynamical extension nevertheless
passes a strong gauge audit: arbitrary time-dependent spherical Weyl
rescalings and diffeomorphisms have exactly zero charge and flux, the entropy
is conformally invariant on the symbolic family, and the normalized
generator extension is unique at linear order under the declared boundary
scalar rule.

\section{Axial Schwarzschild branch splitting}

We now specialize the background to Schwarzschild and use axial
Regge--Wheeler gauge,
\[
 h_{t\phi}=h_0(t,r)S_{\ell m},
 \qquad
 h_{r\phi}=h_1(t,r)S_{\ell m}.
\]
The \(\ell=2\) linearized Bach tensor is fourth order, trace-free, and obeys
the exact linearized Noether identity.  Eliminating the constraint on the
Einstein subbranch reproduces the Regge--Wheeler equation for
\(\Psi=Bh_1/r\),
\begin{equation}
 B(B\Psi')'+(\omega^2-V_{\rm RW})\Psi=0,
 \qquad
 V_{\rm RW}=B\left(\frac6{r^2}-\frac{6m}{r^3}\right).
\end{equation}

\begin{theorem}[Axial composition identity]
On a Ricci-flat background, for the certified axial carrier,
\begin{equation}\label{eq:axial-split}
 \boxed{\delta B_{ab}
 =\frac12\Box(\delta R_{ab})
  +C_{acbd}(\delta R)^{cd}.}
\end{equation}
Consequently the Einstein branch \(\delta R_{ab}=0\) injects into the Bach
kernel, while the extra branch is the second-order system
\begin{equation}
 \frac12\Box\psi_{ab}+C_{acbd}\psi^{cd}=0,
 \qquad \psi_{ab}:=\delta R_{ab}.
\end{equation}
\end{theorem}

The identity is componentwise and off shell in the perturbation.  It does
not extend in this two-term form to the certified non-Einstein static
fixture.  Thus the clean branch split is a Ricci-flat-background theorem.

\section{The extra branch reaches the future horizon}

In ingoing coordinates \((v,r)\), Fourier modes are written
\(e^{i\omega v}\).  The extra carrier's divergence constraint solves its
third component polynomially; the apparent Schwarzschild-coordinate pole is
a chart artifact.

\begin{theorem}[Analytic ingoing family]
For axial \(\ell=2\) and every frequency \(\omega\), the future horizon
\(r=2m\) is a regular singular point of the extra first-order radial system.
Its residue spectrum is
\begin{equation}
 \{0,0,-4im\omega,-2-4im\omega\},
\end{equation}
and the zero eigenspace has dimension two.  Hence there is a two-parameter
family of extra-branch solutions analytic at the future horizon, with no
leading logarithm.
\end{theorem}

Future-horizon regularity therefore does not select the Einstein branch.

\section{Action-derived horizon flux}

Substituting two on-shell Regge--Wheeler solutions into the current
\eqref{eq:green-identity} gives
\begin{equation}\label{eq:rw-null}
 F^r_{EE}
 =-\frac{192\pi\alpha(\omega_1^2-\omega_2^2)}
          {5\omega_1\omega_2r}\Psi_1\Psi_2.
\end{equation}
For conjugate pairs \(\omega_2=\pm\omega_1\), this vanishes identically.

\begin{theorem}[Axial flux disposition]
On Schwarzschild, axial \(\ell=2\):
\begin{enumerate}
\item the Einstein/Regge--Wheeler self-block is symplectically null for
      conjugate wave pairs;
\item ingoing extra modes have a nonzero extra--extra horizon-flux pairing;
\item the Einstein--extra cross pairing is nonzero.
\end{enumerate}
The last two statements are certified at \(m=1\) and independent real
frequencies \(\omega=3/5,2/7,1/2\), using order-16 ingoing series and
radius-stability checks.  The exact Regge--Wheeler null result is the
in-run control.
\end{theorem}

For \(\alpha>0\), the certified ingoing extra fixtures obey
\(iF^r_{XX}>0\) in the declared convention.  We do not promote that classical
sign to a Hilbert norm.  The result needed here is nondegeneracy: the extra
branch is not a zero-flux artifact.

\section{Asymptotic propagation and causal disposition}

At real frequency and large radius, the extra carrier has dispersion
\begin{equation}
 (\lambda_r^2-\omega^2)^2=0
\end{equation}
and growth exponents
\begin{equation}
 \sigma\in\{\pm2im\omega,\ \pm2im\omega-1\}.
\end{equation}
Thus it is luminal, with Coulomb logarithmic phases and amplitudes of order
\(r^0\) or \(r^{-1}\).  There is no growing real-frequency asymptotic
solution.  The Einstein control has the same leading characteristics and
logarithmic phases.

\begin{theorem}[Axial causal truncation no-go]
On the declared Schwarzschild axial \(\ell=2\), real-frequency mode domain,
no boundary prescription based only on future-horizon regularity and local
causal decay at the outer boundary excludes the extra branch while retaining
the Einstein branch.  Selecting the Einstein branch requires an additional
branch projection on scattering data, rather than a local causal
initial-boundary condition.
\end{theorem}

\begin{proof}
The extra branch has a two-dimensional ingoing-analytic horizon family, so
horizon regularity admits it.  Its real-frequency outer solutions are
bounded and have the same characteristic falloff class as the Einstein
solutions, so local asymptotic decay does not distinguish it.  Its self and
mixed Lee--Wald fluxes are nonzero, so it cannot be discarded as a radical
direction.  A condition that removes its independent scattering amplitudes
must therefore project branch data beyond these local endpoint conditions.
\end{proof}

The conclusion is compatible with Euclidean or two-boundary Einstein-branch
selection: such a prescription may define a sector.  The theorem says that
the sector is not generated by ordinary one-ended Lorentzian regularity and
decay in this axial laboratory.

\section{Polar branch: exact split and open causal chain}

For polar \(\ell=2\) Regge--Wheeler-gauge perturbations, the trace variation
\(\delta R\) need not vanish.  The axial formula acquires two universal
trace terms.

\begin{theorem}[Polar trace-coupled composition identity]
On a Ricci-flat background,
\begin{equation}\label{eq:polar-split}
 \boxed{
 \delta B_{ab}=\frac12\Box(\delta R_{ab})
 +C_{acbd}(\delta R)^{cd}
 -\frac16\nabla_a\nabla_b\delta R
 -\frac1{12}g_{ab}\Box\delta R.}
\end{equation}
The Einstein polar branch injects by \(\delta R_{ab}=0\), while the extra
branch is a second-order trace-coupled Lichnerowicz system for
\((\delta R_{ab},\delta R)\).
\end{theorem}

This opens, but does not complete, the polar analogue of the axial chain.
The Zerilli benchmark, horizon indicial system, action-derived flux blocks,
outer asymptotics, and causal disposition remain pending.  No parity-complete
black-hole truncation theorem is claimed until those steps close.

\section{Result ledger}

\begin{center}
\begin{tabularx}{0.98\textwidth}{@{}Y Y Y@{}}
\toprule
Question & Certified answer & Scope\\
\midrule
Static Bach-flat black holes? & exact Laurent-family classification and
three-horizon fixture & local/static\\
Differentiable energy? & residual-basic normalized generator, exact static
first law & static slice; linear spherical gauge audit\\
Einstein axial waves present? & yes; exact Regge--Wheeler injection &
Schwarzschild, axial $\ell=2$\\
Extra axial waves present? & yes; second-order Ricci carrier & same\\
Reach future horizon? & yes; two analytic ingoing amplitudes & mode level\\
Carry current? & nonzero mixed and extra flux; Einstein self-block null &
exact identity plus three frequency fixtures\\
Removed by local causal boundaries? & no & real-frequency axial domain\\
Polar split? & exact trace-coupled identity & causal chain still open\\
Stable/ringing? & not established & complex frequencies open\\
\bottomrule
\end{tabularx}
\end{center}

\section{What is not proved}

The paper does not establish:
\begin{itemize}
\item a complete exterior initial-boundary well-posedness theorem;
\item general \(\ell,m\), the completed polar chain, or non-Einstein
      background branch splitting;
\item complex-frequency poles, mode stability, quasinormal modes, or
      ringdown waveforms;
\item a second-order physical-process first law or nonlinear horizon
      dynamics;
\item asymptotically flat conformal-gravity scattering with a complete
      boundary phase space;
\item Hawking states, particles, a quantum BRST quotient, positivity, or
      unitarity.
\end{itemize}
In particular, ``causally unavoidable'' means unavoidable by the declared
local axial boundary tests, not that every nonlinear or quantum state must
contain the branch.

\section{Certificate and reproduction ledger}

The manuscript is an \textsc{outline-allowed}/\textsc{draft-allowed}
integration of committed certificates at repository baseline
\texttt{35ae48835e4ef0bafb87bb94ff515446879d632b}.

\begin{center}
\begin{tabularx}{0.98\textwidth}{@{}Y Y@{}}
\toprule
Claim & Authoritative certificate\\
\midrule
Static spherical family &
\cert{black_hole_programme/certificates/BH0_STATIC_SPHERICAL_BACKGROUND.json}\\
Normalized generator and first law &
\cert{black_hole_programme/certificates/BH1A_NORMALIZED_GENERATOR.json}\\
Linear spherical dynamical extension &
\cert{black_hole_programme/certificates/BH1B_DYNAMICAL_EXTENSION.json}\\
Axial operator and branch split &
\cert{black_hole_programme/certificates/BH2A_AXIAL_OPERATOR.json}\\
Extra horizon reach &
\cert{black_hole_programme/certificates/BH2A_HORIZON_REACH.json}\\
Action-derived axial current &
\cert{black_hole_programme/certificates/BH2A_FLUX_MATRIX.json}\\
Mixed and extra horizon flux fixtures &
\cert{black_hole_programme/certificates/BH2A_CROSS_FLUX.json}\\
Axial causal disposition &
\cert{black_hole_programme/certificates/BH2A_CAUSAL_DISPOSITION.json}\\
Polar trace-coupled split &
\cert{black_hole_programme/certificates/BH2B_POLAR_SPLIT.json}\\
\bottomrule
\end{tabularx}
\end{center}

Every certificate records a schema, exact or controlled-numerical payload,
independent verifier, mutations or controls, command receipt, dependency
tags, and missing claims.  This manuscript must not advance to
\textsc{theorem-frozen} until its generated claim map and the polar scope are
reviewed.  The axial theorem itself is already certificate-complete.

\section{Outlook}

The next calculation is not a generic stability claim.  It is the polar
analogue of the exact axial sequence:
\[
 \text{Zerilli control}
 \longrightarrow \text{horizon reach}
 \longrightarrow \text{Lee--Wald flux}
 \longrightarrow \text{outer disposition}.
\]
Only then should complex-frequency structure and quasinormal boundary
conditions be activated.  A later asymptotic phase-space construction must
decide which generator is physical, which flux reaches null infinity, and
whether the extra branch survives a globally differentiable charge algebra.

The present result already fixes an important baseline:
\begin{center}
\fbox{\parbox{0.82\textwidth}{\centering
The extra axial branch is regular, radiative, and not removable by ordinary
local causal boundary conditions.}}
\end{center}

\begingroup
\raggedright
\begin{thebibliography}{99}
\bibitem{Bach}
R.~Bach,
\emph{Zur Weylschen Relativit\"atstheorie und der Weylschen Erweiterung des
Kr\"ummungstensorbegriffs}, Math. Z. \textbf{9} (1921) 110--135.

\bibitem{Riegert}
R.~J.~Riegert,
\emph{Birkhoff's theorem in conformal gravity},
Phys. Rev. Lett. \textbf{53} (1984) 315--318.

\bibitem{MannheimKazanas}
P.~D.~Mannheim and D.~Kazanas,
\emph{Exact vacuum solution to conformal Weyl gravity and galactic rotation
curves}, Astrophys. J. \textbf{342} (1989) 635--638.

\bibitem{ReggeWheeler}
T.~Regge and J.~A.~Wheeler,
\emph{Stability of a Schwarzschild singularity},
Phys. Rev. \textbf{108} (1957) 1063--1069.

\bibitem{Zerilli}
F.~J.~Zerilli,
\emph{Effective potential for even-parity Regge--Wheeler gravitational
perturbation equations}, Phys. Rev. Lett. \textbf{24} (1970) 737--738.

\bibitem{LeeWald}
J.~Lee and R.~M.~Wald,
\emph{Local symmetries and constraints},
J. Math. Phys. \textbf{31} (1990) 725--743.

\bibitem{IyerWald}
V.~Iyer and R.~M.~Wald,
\emph{Some properties of Noether charge and a proposal for dynamical black
hole entropy}, Phys. Rev. D \textbf{50} (1994) 846--864,
\href{https://arxiv.org/abs/gr-qc/9403028}{arXiv:gr-qc/9403028}.

\bibitem{Maldacena}
J.~Maldacena,
\emph{Einstein gravity from conformal gravity},
\href{https://arxiv.org/abs/1105.5632}{arXiv:1105.5632}.

\bibitem{Stelle}
K.~S.~Stelle,
\emph{Renormalization of higher-derivative quantum gravity},
Phys. Rev. D \textbf{16} (1977) 953--969.
\end{thebibliography}
\endgroup

\end{document}
